\hypertarget{employee-benefits}{%
\section{Employee benefits}\label{method-corporate-extensions-employee-benefits}}

The additional methods below use the selected IAS 19, IAS 36, IAS 37, IFRS 2, IFRS 13 and IAS 23 provisions in the original 2023 EU incorporating regulation \citep{eu20231803}. This is historical mechanism support; later amendments and actual entity facts remain separate applicability inputs.

DefinedContributionState retains service periods, promised contribution formula, contributions due/paid and applicable legal terms. For earned service requiring contribution 10, debit employee expense 10/credit payable 10; payment settles it. Calling a plan defined contribution needs the limited-obligation facts; a guaranteed benefit or additional funding obligation cannot disappear by name.

The defined-benefit obligation (DBO) measures accrued employee benefits under the stated actuarial assumptions. DefinedBenefitState retains employee/cohort service, benefit formula, projected salary, retirement/turnover/mortality assumptions, discount-curve currency/term, obligation, plan-asset fair value, asset ceiling, contributions, benefits and remeasurements. IAS 19.67--74 requires projected unit credit and the stated service attribution; .83 supplies discount-rate requirements. Assumptions are separately admitted measurements, not constants provided by the accounting rule. A one-employee toy plan promises 100 for each completed service year, paid after year two, with certain completion, zero salary growth and supported 10\% discount. At year-one end accrued benefit 100 due in one year has PV 1000/11. In year two that opening obligation accretes 100/11 and the second earned benefit adds 100; obligation 200 then settles 200. Merely using current salary times years without timing would give the wrong year-one obligation.

A separate funded-plan roll-forward supplies an independent statement example: opening DBO 100, plan assets 80 and net liability 20; current service 12; interest 5 on DBO, interest 4 on assets at 5\%; contributions 10 at year end; benefits 8 paid from plan at year end; actuarial obligation loss 7; actual asset return 6. DBO closes 100+12+5\ensuremath{-}8+7=116, assets 80+6+10\ensuremath{-}8=88, net liability 28. Expense is 12+1=13, OCI loss is 7\ensuremath{-}(6\ensuremath{-}4)=5. Debit service expense 12/credit net pension liability 12; debit interest expense 1/credit liability 1; debit OCI 5/credit liability 5; debit liability 10/credit cash 10. End liability 20+13+5\ensuremath{-}10=28. Plan-paid benefits reduce both DBO and plan assets and do not cause a second employer cash payment. IAS 19.113, .120--127 supports these components; remeasurement OCI is not recycled to profit under 122. A surplus requires the asset ceiling and applicable minimum funding analysis, not recognition of any calculated negative liability.

Unknown actuarial inputs give conditional obligations or measurement\_unresolved. Curtailment, settlement, amendment and multi-employer/insured arrangements need their own fact conditions and measurement-date remeasurement. A buyout is not settlement if obligations remain. Sensitivities vary one assumption only when that conditional meaning is explicit; correlated risks require joint scenarios.

\hypertarget{impairment-and-provisions}{%
\section{Impairment and provisions}\label{method-corporate-extensions-impairment-and-provisions}}

ImpairmentState retains asset/CGU boundaries, carrying amounts, allocated goodwill, corporate assets, recoverable-amount inputs and previous impairments/depreciation. IAS 36.9--10 determines indicators and mandatory annual tests. Recoverable amount is the higher of value in use and fair value less disposal costs (.18); it is not the lower of two pessimistic scenarios. VIU needs the supported cash flows, forecast horizon/growth, discount and exclusion of inappropriate financing/tax or uncommitted restructuring effects under 33--57. Compute conditional values if those estimates are unknown; do not infer them from historical profitability.

For a cash-generating unit (CGU), assets excluding goodwill 100 plus goodwill 20, recoverable amount 90: loss 30 first removes goodwill 20 and then reduces other assets 10 subject to their floors under 104--105. Debit impairment 30/credit goodwill 20/credit asset allowance 10. If the asset floor is 95 the requested reduction to 90 is impermissible; apply the remaining allocation/recognition conditions, not a forced plug. For a standalone cost-model asset with carrying 80, recoverable 60, recognize 20 loss. At a later measurement before further depreciation, supported recovery to 90 permits reversal only 20 to the no-impairment 80 ceiling (.117--119). Goodwill impairment is not reversed (.124). Different dates require the hypothetical unimpaired depreciation path, not the old nominal cost as cap.

\subsection{An asset floor changes the allocation of a unit's loss}
\label{sec:impairment-floors}

A cash-generating unit can lose value while one of its assets retains a
substantial disposal value. Allocating the unit's loss in proportion to asset
balances just once may then reduce that asset too far. IAS 36.104--105 first
allocates loss to goodwill, then to the other assets pro rata, while protecting
each asset from a reduction below the highest of zero, its fair value less
costs of disposal if measurable, and its value in use if determinable. A share
that cannot be assigned to one asset is redistributed to the others
\citep{eu20231803}.

Consider a synthetic cost-model unit containing equipment A with carrying
amount 60, equipment B with carrying amount 40, and allocated acquired
goodwill 20. Cash 80 lies outside the tested unit. With no liabilities, opening
equity is 200. The unit boundary and goodwill allocation are assumed to have
been admitted under paragraphs 66--80; there are no shared corporate assets,
noncontrolling interests or liabilities requiring an adjustment to this
comparison. In a real case those judgments precede the calculation. In
particular, paragraph 78 requires consistent treatment of a liability that is
included in the unit's recoverable value; subtracting it from only the carrying
amount would compare different populations.

At measurement date $T_0$, admitted value in use is 70 and fair value less
disposal costs is 65, so the unit's recoverable amount is 70. Individual
values in use are supported as not determinable, while the individual disposal
values are 55 for A and 5 for B. These are supplied valuation assumptions,
not values inferred by an allocation routine. They establish floors of 55
and 5. A missing valuation would be a different state: the program may omit
a term because its non-determinability is admitted, but may not replace an
unexamined input with zero.

The unit's loss is $60+40+20-70=50$. Goodwill takes the first 20. An
unconstrained allocation of the remaining 30 would assign
$30\times60/100=18$ to A and $30\times40/100=12$ to B. A, however, can
absorb only $60-55=5$. Its excess share of 13 passes to B, whose total loss
becomes $12+13=25$, within its capacity $40-5=35$. The closing balances are
therefore
\[
  A=60-5=55,\qquad B=40-25=15,\qquad
  \text{goodwill}=20-20=0.
\]
Together they equal the unit's recoverable amount of 70. Debit impairment
expense 50, credit goodwill 20, and credit the separate asset impairment
allowances by 5 and 25. Cash 80 plus equipment 70 equals equity 150. The
unconstrained 18/12 entry would also balance, but would leave A at 42,
below its admitted floor of 55. Clipping A's loss to 5 without redistributing
the excess would instead recognize only 17 of the required remaining 30.

For an implementer, each round has a remaining amount and an active set of
assets. Compute shares from their pre-allocation carrying amounts. Fix each
share that reaches an asset's available capacity, remove that asset from the
active set, and redistribute the remaining amount. A round that needs further
redistribution removes at least one asset, so the procedure terminates after
finitely many such removals. For positive carrying weights $b_i$ and
capacities $c_i$, its allocation has the form
\[
 x_i=\min(c_i,\lambda b_i),\qquad
 \sum_i x_i=\min\!\left(q,\sum_i c_i\right),
\]
where $q$ is the requested allocation and $\lambda\geq0$ gives the common
proportional rate among assets that have not reached their limits. In the
example A is fixed at 5 and B receives the remaining 25. A source floor above
an asset's existing cost gives zero reduction capacity; it does not create
permission to increase that asset during an impairment calculation.

The capacity limit also explains how to handle an unsuccessful allocation.
A diagnostic input requesting 50 against capacities 5 and 35 can allocate
only 40, leaving 10. That arithmetic does not establish that the supplied
unit and individual valuations are jointly coherent. It calls for checking
measurement assumptions, the unit boundary and any separate obligations.
Paragraph 108 permits a liability for a remaining unit loss only if another
IFRS requires one. AccDSL must retain the residual and the unresolved decision;
it cannot invent a payable to make the unit reach a requested total. This
diagnostic is not a second admitted company example.

\subsection{A recovery is limited by the depreciation that would have occurred}
\label{sec:impairment-reversal-memory}

Suppose both pieces of equipment have five remaining years of even use and
zero residual values at $T_0$. Under paragraph 63, their revised carrying
amounts give next-year depreciation $55/5=11$ and $15/5=3$. At $T_1$, one
full year later, they carry 44 and 12. The company must also retain the path
that would have occurred without impairment. On that path depreciation would
have been $60/5=12$ and $40/5=8$, leaving 48 and 32. Restoring a past loss
therefore cannot mean restoring the old nominal balances.

Assume a supported change in expected future cash flows now gives the unit a
recoverable amount of 90. Individual recoverable amounts remain admitted as
not determinable. Paragraphs 114--116 require the change in estimates; mere
passage of time and discount unwinding would not support this reversal.
Paragraphs 122--123 allocate a unit reversal over assets other than goodwill
pro rata, subject to each asset's recoverable amount if determinable and its
depreciated no-impairment amount. Goodwill itself cannot be restored under
paragraph 124 \citep{eu20231803}.

The requested increase is $90-(44+12)=34$. A has room for only $48-44=4$;
B has room for $32-12=20$. The same redistribution procedure therefore
recognises $4+20=24$, leaving equipment at 48 and 32 and goodwill at zero.
The recognized total of 80 is below the unit's value of 90. That difference
does not identify an omitted asset, authorize goodwill recognition or imply
cash has been received.

Debit the impairment allowances by 4 and 20 and credit reversal gain 24.
The original allowances of 5 and 25 consequently retain balances of 1 and 5.
Accumulated depreciation is 11 and 3. The register now reconciles as
\[
  \text{PPE}=100-\underbrace{(11+3)}_{\text{depreciation}}
                 -\underbrace{(1+5)}_{\text{impairment allowances}}=80.
\]
The remaining allowances record the different past depreciation paths.
Erasing all 30 of the original allowances would produce PPE of 86, exceeding
the no-impairment ceiling by 6. Future annual depreciation is $48/4=12$ and
$32/4=8$, consistent with the remaining lives.

\begin{center}
\begin{tabular}{@{}lrr@{}}
\toprule
Pretax amount & $T_0$ impairment & Following year to $T_1$\\
\midrule
Impairment expense & 50 & 0\\
Depreciation / reversal gain & 0 / 0 & 14 / 24\\
Profit / OCI & $-50$ / 0 & 10 / 0\\
Closing cash / PPE & 80 / 70 & 80 / 80\\
Closing goodwill / liabilities & 0 / 0 & 0 / 0\\
Closing equity & 150 & 160\\
Operating cash & 0 & 0\\
\bottomrule
\end{tabular}
\end{center}

The first column isolates the impairment measurement; the second follows the
next full year. With no cash or owner movements in this slice, the operating
cash reconciliations are $-50+50=0$ and $10+14-24=0$, removing the noncash
loss, depreciation and reversal under IAS 7.20. Equity moves from 200 to
150, then to 160. Tax remains a separate calculation under IAS 36.64 and
IAS 12. These technical rules use the retained original 2023 EU text, with
its date corrigendum; the relative dates do not assert current reporting
admission or a complete later amendment review.

The method thus needs actual and no-impairment depreciation histories,
measurement dates, individual valuation-status judgments, goodwill identity
and evidence of the changed estimate. There is also an explicit numerical
boundary. If every eligible asset has zero carrying amount but a positive
reversal is requested, the positive-weight procedure above has no denominator.
It returns an unresolved allocation basis and preserves the history; it does
not choose equal weights without an additional justified interpretation.
That extension, along with shared corporate assets and other scope
interactions, remains distinct from the fully specified positive-weight case.


\subsection{Recognising a separate obligation}

ProvisionState retains obligating past event, present legal/constructive duty, outflow assessment, estimate distribution, settlement dates, risk/discount basis, reimbursements and evidence. IAS 37.14 requires present obligation, probable outflow and reliable estimate; .36--47 governs best estimate/risk/discount. An intended future operating expenditure without that duty is not a provision. Assume an established obligation to pay 121 in two years, deterministic amount, appropriate 10\% rate and no double-counted risk: recognize expense/provision 100, accrete interest 10 then 11, settle 121. The passage-of-time increase is borrowing cost under 60. If the obligation is an eligible asset restoration cost, initial debit follows that asset's method rather than always expense; later changes need the applicable restoration interpretation. Contingent liabilities/disclosures remain visible when recognition conditions are not met; no generic 0.5 threshold may substitute for the relevant standard's required analysis.

\hypertarget{share-based-compensation}{%
\section{Share-based compensation}\label{method-corporate-extensions-share-based-compensation}}

AwardState distinguishes equity versus cash settlement, grant/vesting/service dates, granted units, grant-date fair value, service/nonmarket/market conditions, expected units vesting, cumulative expense and changes/cancellations. IFRS 2.7--11, 15--23 supports the employee equity-settled method. With 100 units, grant value 3, two-year service condition and all expected to vest, year-one cumulative expense 150 and equity 150. At vesting only 80 satisfy the service condition: cumulative 240, second-year expense 90. Debit employee expense and credit equity. Failure of a market condition is different under 21; if the employee satisfies other vesting conditions, it does not automatically reverse the grant-date measured service expense. Unknown conditions cannot be represented by a single forfeiture percentage without their type and evidence.

Cash-settled awards remeasure the liability under 30--33. Same 100 expected units, half service and current unit fair value 4 gives liability/expense 200 at year one; 80 actual units and value 5 at settlement gives total 400, second expense 200, then debit liability 400/credit cash 400. Reusing the equity grant value 3 would be wrong. Actual award valuation requires the specified model/market inputs, option rights, volatility, dividends and calibration; the fair values 3/4/5 are case inputs. Modifications and settlement-choice arrangements require their own conditions, with the original minimum expense and increment analyzed under the applicable provisions rather than overwriting the original grant.

\hypertarget{fair-value-and-capitalization-boundaries}{%
\section{Fair value and capitalization boundaries}\label{method-corporate-extensions-fair-value-and-capitalization-boundaries}}

A ValuationRequest names the unit of account, market/date, asset/liability characteristics, principal/accessible market, technique, input hierarchy, counterparty/own non-performance risk where relevant, calibration and sensitivity. IFRS 13.24, .42--43 and .67--80 guide the selected measurement. Quoted identical asset price 7 for 10 units gives 70 in a qualifying Level 1 case; a blockage discount based only on holding size cannot reduce it to 63 under 69/80. A present-value illustration 111 due next year at an independently supported 11\% rate gives 100; that equality does not establish the rate, market comparability or accounting classification. Historical EIR, an impairment discount, risk-neutral pricing and a physical expected return answer different questions.

Borrowing-cost capitalization requires a qualifying asset, eligible costs, commencement/suspension/cessation dates and the specific/general-borrowing calculation, before moving interest into cost. IAS 23.8--25 supplies these rules in the retained historical EU text: for specific borrowing, eligible actual costs less temporary investment income; for general borrowing, the applicable weighted rate and cost ceiling. Commencement requires expenditures, borrowing costs and necessary preparation activities; extended inactive periods can suspend it, and substantial completion ends it. Assume specific borrowing 100 outstanding during a full eligible construction year at 5\%, with an admitted expenditure schedule and temporary investment income 1 on temporarily unused funds during that period. Net capitalized cost is 4. Real admission still needs the dated standard, actual schedule and eligibility evidence. It is not a default 5\% add-on to all assets. Internally developed intangibles likewise need research/development separation and all applicable recognition criteria; a forecast of future sales does not establish an asset.

The extension interface returns the stated measured components, proposed journal, economic delta, disclosures and unresolved inputs under the event-basis, authorization and commit contracts. Calculation success does not permit another operation to fill a missing authority/fact with a default. Real-company acquisition/tax/employee/valuation evidence and the full exceptions remain explicit internal preparation, with proposed answers and competing treatments retained for any later authorized domain review.
